\scanpage{105}
% archive foliation: 104
\fbox{\parbox{.95\linewidth}{\uline{Poids de la cohomologie \inserted{$l$-adique} d'une variété singulière ou/et non propre}. \inserted{(sur corps $k$ \unsure{p. f. $0$})}}}

\textcircled{A} \uline{Coefficients constants}

\textcircled{1} Soit d'abord $X$ lisse et \unsure{quasi-projective|quasi-projectif}, donc par résolution forte (que nous admettons).--- on aura
\[
X=X'-Y\qquad Y=\bigcup_iY_i
\]
$X'$ lisse et projective, $Y_i$ diviseurs lisses, \inserted{irred.} à croisements normaux mutuels. \deleted{\unsure{Si}} On a donc
\diagram{\begin{tikzcd}[column sep=small]
\inserted{H^{i-1}(X')} \arrow[r] & \substack{\inserted{H^{i-1}(Y)}\\\deleted{\unsure{R^if_!}}} \arrow[r] & H^i_!(X) \arrow[r,dash] & H^i(X') \arrow[r,dash] & H^i(Y) \arrow[r,dash] & H^{i+1}_!(X) \arrow[r] & \cdots
\end{tikzcd}}

\margin{\diagram{\begin{tikzcd}
X \arrow[r,phantom,"\subset"] \arrow[d,"f"'] & X' \arrow[dl,"f'"]\\
\operatorname{Sp}k &
\end{tikzcd}}}

\deleted{\unsure{donc}} donc les poids dans $H^i_!(X)$ [qui se déterminent via val. absolues de val. propres de Frob, par réduction au cas d'un corps de type fini \unsure{\ldots}) \unsure{sont contenus} dans ceux de $H^i(X')$ [\unsure{réduits} à $i$ par Weil] et ceux de $H^{i-1}(Y)$.

On a
\[
H^{\unsure{n}}(Y)\Longleftarrow E^{pq}_2
=H^p(i_0\ldots i_r\longmapsto H^q(Y_{i_0}\cap\cdots\cap Y_{i_r}))
\]
\inserted{on peut se borner ici aux cochaines \uline{alternées}}
\margin{$\left\{\begin{array}{l}
\text{de poids }q\\
\text{de niveau }\le q\\
[\text{intéressant si}\\
i\unsure{=|\equiv}p+q\le n-1]\\
\text{de niveau }\le\unsure{2n-i-q}\ \deleted{\unsure{q}}\\
[\text{si }i=p+q\ge n]
\end{array}\right.$}
\margin{$(2n-i)-q$}

\deleted{Donc} Or ici $Y_{i_0}\cap\cdots\cap Y_{i_r}$ est projectif et lisse (de dimension $n-r-1$ si $X$ \deleted{\unsure{est}} purement de dim $n$). Donc le poids des termes $E^{pq}_2$ est $q$, et on trouve que le poids de $H^j(Y)$ est $\le j$. Donc :

\deleted{\unsure{Proposition}} \uline{Proposition 1} $H^i_!(X)$ est de poids \inserted{entre $0$ et} $\le i$. La partie de poids $i$ (resp. de poids $\le i-1$) est le quotient image de $H^i_!(X)\mathbin{-}H^i(X')$ [resp. l'image de $H^{i-1}(Y)\longrightarrow H^i_!(X)$). De plus, si $i\ge n=\dim X$, alors le \uline{niveau} de $H^i_!(X)$ est $\le2n-i$, et les deux \unsure{derniers} morceaux précédents sont de niveaux $\le2n-i-1$; de plus, il n'y a de poids qu'entre $2(i-n)$ [car il faut écrire $q\le2(n-p-1)$ !] et $i$, i.e \unsure{un gr. pondéré} d'amplitude $2n-i$ \ldots

\margin{\fbox{\parbox{.24\linewidth}{\uline{NB} Tous les morceaux de la filtration sont ``effectifs''}}}
\margin{\fbox{\parbox{.24\linewidth}{En tous cas niveau $\le\operatorname{Inf}(i,n)$}}}

Pour ce complément, notons que le niveau de $E^{pq}_2$ est $\le$ celui de $H^q(\underbrace{Y_{i_0}\cap\cdots\cap Y_{i_p}}_{\dim n-p-1})$, donc est
\[
\le \deleted{\unsure{2n-p-q}}\,
2(\deleted{\unsure{n}}n-p-1)-q
=\deleted{\unsure{2n-(p+q)}}(2n-i-2)-p
\inserted{=2(n-1)-i-p}
\quad(i=p+q)\quad\text{donc }\le2(n-1)-i
\]
